\documentclass[11pt,a4paper]{article}
\usepackage{fontspec,amsmath,amssymb,amsthm,microtype}
\setmainfont{lmroman10-regular.otf}[BoldFont=lmroman10-bold.otf,ItalicFont=lmroman10-italic.otf,BoldItalicFont=lmroman10-bolditalic.otf]
\usepackage{xcolor}
\usepackage[most]{tcolorbox}
\usepackage{mathrsfs}
\providecommand{\lt}{<}
\providecommand{\gt}{>}
\setmonofont{DejaVuSansMono.ttf}[BoldFont=DejaVuSansMono-Bold.ttf,ItalicFont=DejaVuSansMono-Oblique.ttf,BoldItalicFont=DejaVuSansMono-BoldOblique.ttf,Scale=0.86]
\newfontfamily\leanfallback{STIXTwoMath-Regular.otf}[Scale=0.86]
\definecolor{leanrule}{gray}{0.62}
\definecolor{leanback}{gray}{0.962}
\definecolor{leantitle}{gray}{0.22}
\definecolor{leanmuted}{gray}{0.38}
\newtcolorbox{leanbox}[1]{enhanced,breakable,colback=leanback,colframe=leanrule,
  boxrule=0.4pt,arc=0.6mm,left=2.4mm,right=2.4mm,top=1.8mm,bottom=1.8mm,
  before skip=0.8em,after skip=1em,fonttitle=\small\bfseries,coltitle=leantitle,
  colbacktitle=leanback,titlerule=0.3pt,title={#1}}
\newlength{\leancharwidth}
\newenvironment{leancode}{\par\smallskip\begingroup\ttfamily\footnotesize\raggedright
  \settowidth{\leancharwidth}{M}\setlength{\parskip}{0pt}\setlength{\parindent}{0pt}
  \frenchspacing}{\par\endgroup\smallskip}
\newcommand{\leanline}[2]{\hangindent=\dimexpr#1\leancharwidth+4\leancharwidth\relax
  \hangafter=1\noindent\hspace*{#1\leancharwidth}#2\par}
\newcommand{\leansource}[1]{\par{\footnotesize\color{leanmuted}#1}\par}
\newcommand{\leanlabel}[1]{\par\smallskip\noindent{\small\bfseries #1}\par}
\usepackage[margin=25mm]{geometry}
\usepackage[unicode,hidelinks]{hyperref}
\hypersetup{pdfauthor={LangevinLean Project}}
\setlength{\emergencystretch}{3em}
\setlength{\parskip}{0.4em}
\setlength{\parindent}{0pt}
\allowdisplaybreaks[2]
\raggedbottom
\title{\texorpdfstring{Gibbs invariance of the exact split transition}{Gibbs invariance of the exact split transition}}
\author{LangevinLean Project}
\date{}
\begin{document}\maketitle
Draft mathematical exposition; these notes may contain mistakes.


To measure OBABO discretization bias, keep both thermostat noises exactly the same and replace its Verlet middle portion by the proved exact Hamiltonian flow. This comparison step preserves the Gibbs target. Its invariance and its common-source \(L^{2}\) random variables are constructed here, so the numerical error bound can compare random variables on the same probability space.



Replacing only the central Verlet stage by exact Hamiltonian evolution creates a target-preserving comparison transition. Retaining the same two thermostat noises yields the natural coupling for its local numerical discrepancy.


\section{\texorpdfstring{The exact comparison map with the same two noises}{The exact comparison map with the same two noises}}


Starting from \(z=(v,x)\), apply \(O_{\delta,\sigma,\xi}\), evolve the resulting phase under the selected Hamiltonian flow for time \(h\), and then apply \(O_{\delta,\sigma,\xi\prime}\). The map uses the same first and second innovations as the OBABO update. Only the deterministic Hamiltonian evolution changes. Continuity of the selected Hamiltonian endpoint and of the two affine thermostat maps proves joint continuity in the input phase and the paired noise.



The comparison replaces the central Verlet stage by exact Hamiltonian evolution while retaining the same thermostat half steps. Since all three exact stages preserve the Gibbs target, their composition is a natural target-preserving comparator for the numerical OBABO map.

\begin{leanbox}{Lean: \texttt{exactObaboStep}, \texttt{exactObaboStep\_\allowbreak{}continuous}}
\leansource{\texttt{exactObaboStep} \quad KineticLangevin/\allowbreak{}ExactOBABO.lean, lines 15–18}
\begin{leancode}
\leanline{0}{def\ exactObaboStep\ (f\ :\ E\ →\ ℝ)\ (M\ :\ ℝ≥0)\ (hLip\ :\ LipschitzWith\ M\ (gradient\ f))}
\leanline{4}{(h\ :\ ℝ)\ (hh\ :\ 0≤h)\ (δ\ σ\ :\ ℝ)\ (z\ :\ E×E)\ (ξ\ :\ E×E)\ :\ E×E\ :=}
\leanline{2}{ornsteinStep\ δ\ σ\ ξ.2}
\leanline{4}{(hamiltonianFlow\ f\ M\ hLip\ h\ hh\ (ornsteinStep\ δ\ σ\ ξ.1\ z)\ h)}
\end{leancode}
\leansource{\texttt{exactObaboStep\_\allowbreak{}continuous} \quad KineticLangevin/\allowbreak{}ExactOBABO.lean, lines 20–25}
\begin{leancode}
\leanline{0}{theorem\ exactObaboStep\_continuous\ (f\ :\ E\ →\ ℝ)\ (M\ :\ ℝ≥0)}
\leanline{4}{(hLip\ :\ LipschitzWith\ M\ (gradient\ f))\ (h\ :\ ℝ)\ (hh\ :\ 0≤h)\ (δ\ σ\ :\ ℝ)\ :}
\leanline{4}{Continuous\ (Function.uncurry\ (exactObaboStep\ f\ M\ hLip\ h\ hh\ δ\ σ))}
\end{leancode}
\leanlabel{How the Lean proof follows the argument}
\texttt{exact\allowbreak{}Obabo\allowbreak{}Step} names the mathematical object defined above. Its defining expression appears in the adjacent statement; subsequent identities may unfold this definition directly.




\texttt{exact\allowbreak{}Obabo\allowbreak{}Step\_\allowbreak{}continuous} gives the formal counterpart of this step. The main intermediate facts are \texttt{hc} in that order.
\end{leanbox}

\section{\texorpdfstring{Preserve the first updated phase and its independence of fresh noise}{Preserve the first updated phase and its independence of fresh noise}}


The source law is explicitly \(\pi\otimes(\mathcal N(0,I)\otimes\mathcal N(0,I))\), where \(\pi\) is the normalized Gaussian-times-Gibbs phase target. Reassociate it as \((\pi\otimes\mathcal N(0,I))\otimes\mathcal N(0,I)\). The first thermostat preserves \(\pi\) because \(\delta^2+\sigma^2=1\), while the untouched second Gaussian remains independent. This follows using the measure-preserving product associator and the product pushforward identity. The conclusion is the joint law of the updated phase and the second fresh noise, not just a marginal invariance statement.



After the first thermostat update the phase still has the Gibbs law. The second thermostat innovation remains independent because it was chosen independently of the initial phase and the first innovation. This justifies applying its invariance property after the Hamiltonian stage.

\begin{leanbox}{Lean: \texttt{firstThermostat\_\allowbreak{}joint\_\allowbreak{}fresh}}
\leansource{\texttt{firstThermostat\_\allowbreak{}joint\_\allowbreak{}fresh} \quad KineticLangevin/\allowbreak{}ExactOBABO.lean, lines 29–50}
\begin{leancode}
\leanline{0}{theorem\ firstThermostat\_joint\_fresh}
\leanline{4}{(f\ :\ E\ →\ ℝ)\ (m\ :\ ℝ)\ (hm\ :\ 0<m)\ (hf\ :\ Differentiable\ ℝ\ f)}
\leanline{4}{(hsc\ :\ StrongConvexOn\ univ\ m\ f)\ (δ\ σ\ :\ ℝ)\ (hvariance\ :\ δ\textasciicircum{}2+σ\textasciicircum{}2=1)\ :}
\leanline{4}{((kineticGibbsTarget\ f).prod\ ((stdGaussian\ E).prod\ (stdGaussian\ E))).map}
\leanline{6}{(fun\ p\ ↦\ (ornsteinStep\ δ\ σ\ p.2.1\ p.1,p.2.2))\ =}
\leanline{8}{(kineticGibbsTarget\ f).prod\ (stdGaussian\ E)}
\end{leancode}
\leanlabel{How the Lean proof follows the argument}
\texttt{first\allowbreak{}Thermostat\_\allowbreak{}joint\_\allowbreak{}fresh} gives the formal counterpart of this step. The main intermediate facts are \texttt{\ensuremath{\pi}}, \texttt{G}, \texttt{h\allowbreak{}OU}, \texttt{hfirst}, \texttt{h\allowbreak{}Assoc}, \texttt{hp} in that order. Named proof applications include \texttt{hp.\allowbreak{}symm}; the adjacent proof shows their hypotheses and order of use.
\end{leanbox}

\section{\texorpdfstring{Compose the three proved invariant transitions}{Compose the three proved invariant transitions}}


After the first \(O\) step the joint law is \(\pi\otimes\mathcal N(0,I)\). The exact Hamiltonian map preserves \(\pi\) under the original strong-convexity and Lipschitz-gradient assumptions; applying it only to the phase therefore leaves the same joint law with the second noise. The last \(O\) step preserves \(\pi\) again. Three explicit pushforward composition identities prove that the complete \(O\)–Hamiltonian–\(O\) random map preserves the normalized target. No invariance hypothesis is supplied to this theorem.



Push the target successively through the first thermostat, Hamiltonian flow, and final thermostat. Each pushforward equals the target, so their composition does as well. This proves invariance for the comparison transition, including both noises.

\begin{leanbox}{Lean: \texttt{exactObaboStep\_\allowbreak{}preserves\_\allowbreak{}target}}
\leansource{\texttt{exactObaboStep\_\allowbreak{}preserves\_\allowbreak{}target} \quad KineticLangevin/\allowbreak{}ExactOBABO.lean, lines 54–85}
\begin{leancode}
\leanline{0}{theorem\ exactObaboStep\_preserves\_target}
\leanline{4}{(f\ :\ E\ →\ ℝ)\ (m\ :\ ℝ)\ (hm\ :\ 0<m)\ (M\ :\ ℝ≥0)}
\leanline{4}{(hf\ :\ Differentiable\ ℝ\ f)\ (hsc\ :\ StrongConvexOn\ univ\ m\ f)}
\leanline{4}{(hLip\ :\ LipschitzWith\ M\ (gradient\ f))\ (h\ :\ ℝ)\ (hh\ :\ 0≤h)}
\leanline{4}{(δ\ σ\ :\ ℝ)\ (hvariance\ :\ δ\textasciicircum{}2+σ\textasciicircum{}2=1)\ :}
\leanline{4}{((kineticGibbsTarget\ f).prod\ ((stdGaussian\ E).prod\ (stdGaussian\ E))).map}
\leanline{6}{(Function.uncurry\ (exactObaboStep\ f\ M\ hLip\ h\ hh\ δ\ σ))\ =\ kineticGibbsTarget\ f}
\end{leancode}
\leanlabel{How the Lean proof follows the argument}
\texttt{exact\allowbreak{}Obabo\allowbreak{}Step\_\allowbreak{}preserves\_\allowbreak{}target} gives the formal counterpart of this step. The main intermediate facts are \texttt{\ensuremath{\pi}}, \texttt{G}, \texttt{A}, \texttt{B}, \texttt{C}, \texttt{h\allowbreak{}A} in that order. Named proof applications include \texttt{he.\allowbreak{}symm}; the adjacent proof shows their hypotheses and order of use.
\end{leanbox}

\section{\texorpdfstring{The \(L^{2}\) variables used by the local comparison}{The L² variables used by the local comparison}}


Pull the initial velocity and position through the first product projection, and the two Gaussian innovations through their separate projections. The target projection moment theorems and Gaussian finite moments prove all four coordinates belong to \(L^{2}\). Their \(L^{2}\) representatives agree almost everywhere with the original coordinates. The proved pointwise representation of a lifted thermostat identifies the first updated \(L^{2}\) phase with the first \(O\) map. Projecting the established phase/\allowbreak{}fresh-noise joint law now gives its exact target marginal. These random variables retain the common noise coupling and have the stationary laws needed by the local Hamiltonian–Verlet error theorem.



The initial phase and both innovations are represented jointly on their product space. The exact Hamiltonian comparator is square integrable under its stationary law. These variables provide the common-noise coupling needed to integrate the Verlet local state error.

\begin{leanbox}{Lean: \texttt{exactObabo\_\allowbreak{}lp\_\allowbreak{}source}}
\leansource{\texttt{exactObabo\_\allowbreak{}lp\_\allowbreak{}source} \quad KineticLangevin/\allowbreak{}ExactOBABO.lean, lines 89–127}
\begin{leancode}
\leanline{0}{theorem\ exactObabo\_lp\_source}
\leanline{4}{(f\ :\ E\ →\ ℝ)\ (m\ :\ ℝ)\ (hm\ :\ 0<m)\ (hf\ :\ Differentiable\ ℝ\ f)}
\leanline{4}{(hsc\ :\ StrongConvexOn\ univ\ m\ f)\ (δ\ σ\ :\ ℝ)\ (hvariance\ :\ δ\textasciicircum{}2+σ\textasciicircum{}2=1)\ :}
\leanline{4}{let\ P\ :=\ (kineticGibbsTarget\ f).prod\ ((stdGaussian\ E).prod\ (stdGaussian\ E))}
\leanline{4}{∃\ (Z\ :\ Lp\ E\ 2\ P\ ×\ Lp\ E\ 2\ P)\ (Ξ\ Ξ'\ :\ Lp\ E\ 2\ P),}
\leanline{6}{lpPhaseCoe\ P\ Z\ =ᵐ[P]\ Prod.fst\ ∧}
\leanline{6}{Ξ\ =ᵐ[P]\ (fun\ p\ :\ (E×E)×(E×E)\ ↦\ p.2.1)\ ∧}
\leanline{6}{Ξ'\ =ᵐ[P]\ (fun\ p\ :\ (E×E)×(E×E)\ ↦\ p.2.2)\ ∧}
\leanline{6}{P.map\ (lpPhaseCoe\ P\ (ornsteinStep\ δ\ σ\ Ξ\ Z))=kineticGibbsTarget\ f}
\end{leancode}
\leanlabel{How the Lean proof follows the argument}
\texttt{exact\allowbreak{}Obabo\_\allowbreak{}lp\_\allowbreak{}source} gives the formal counterpart of this step. The main intermediate facts are \texttt{P}, \texttt{V}, \texttt{X}, \texttt{W}, \texttt{hv}, \texttt{hx} in that order. Named proof applications include \texttt{Prod.\allowbreak{}ext}, \texttt{hj}; the adjacent proof shows their hypotheses and order of use.
\end{leanbox}

\end{document}
